% ===============================================
% LaTeX Task Template
%
% Author: Matthias Nickel
% E-Mail: Matthias Nickel@tu-dresden.de
% Kudos to Julian Haase (Julian.Haase@tu-dresden.de) for creating the original template.
% Date: 02.10.2025
% Please contact the author in case of questions.
% ===============================================
\RequirePackage[ngerman=ngerman-x-latest]{hyphsubst}	% better german hyphenation rules
\documentclass[english, 10pt]{tudscrreprt}		% load Koma-Script-Class for the CD of the TUD
														% Replace 'english' with 'ngerman' for the German version and vise versa.
\usepackage{scrhack}									% includes koma script patches for other packages
% ===============================================
%	Character coding and fonts for the document
% ===============================================
\usepackage[utf8]{inputenc}	% UTF-8 Support
\usepackage[T1]{fontenc}	% T1 Fonts for font encoding
\usepackage{babel} %
\usepackage{microtype}	% better layout of the text
\usepackage[sfdefault]{noto}
% ===============================================
%	Watermark - only for the draft
% ===============================================
% \usepackage[firstpage]{draftwatermark}
% \SetWatermarkText{\textbf{draft}}
% \SetWatermarkScale{1.1}
% ===============================================
%	Additional packages
% ===============================================
\usepackage{iflang} 						% For correct headings
\usepackage{isodate}						% for several date formats
\usepackage{enumitem}\setlist{noitemsep}	% less space between the items
\usepackage{tudscrsupervisor}				% package for supervising templates for the CD of the TUD

% ===============================================
%	Wrapper for discipline
% ===============================================
\newcommand{\module}[1]{\discipline{#1}}	% Wrapper to use 'module' instead of 'discipline'
\TUDoptions{extrabottommargin=-5cm}
%
\begin{document}
% ===============================================
%	Heading for the chair
% ===============================================
\IfLanguageName{english}{% english
	\faculty{Faculty of Computer Science}
	\institute{Institute of Computer Engineering}
	\chair{Chair of Adaptive Dynamic Systems}
}{% Deutsch
	\faculty{Fakultät Informatik}
	\institute{Institut für Technische Informatik}
	\chair{Professur für Adaptive Dynamische Systeme}
}%
\headlogo{logo/ADS_Logo.png}				% Comment out this line if the ADS logo should be used.
% ===============================================
% ===============================================
% ===============================================
%	Change this for the given case
% ===============================================
% ===============================================
% ===============================================
\title{%
	FPGA-based Real-Time Learning Architecture with Partial Reconfiguration for Spiking Neural Networks
}
\thesis{master} % master, diploma, bachelor, student, project, evidence, internship or seminar -> for details take a look at table 2.1 in tudscr.pdf
\author{
	Yike Jing
	\matriculationnumber{5201619}%
	% \module{Master-Thesis-NES}%											% Module ID, e.g., INF-PM-ANW, Eul-NES-C-PW; Comment out for theses
	\course{Master NES}%								% Course of studies, e.g., Master Computer Science, Diploma Information Systems Engineering, Master Nanoelectronic Systems
}
\matriculationyear{2023}
\issuedate{15.05.2026}
\duedate{23.10.2026}
\supervisor{Gloria Sepanta \and Dr.-Ing. Lester Kalms }											% If multiple supervisors use '\and' to separate them.
\referee{(TU Dresden) Prof. Dr.-Ing. Diana Göhringer \and (TU Dresden) Prof. Dr. rer. nat. Josef Weidendorfer \and (KU Leuven) Prof. Dr. Ir. Nele Mentens}		% If multiple reviewers use '\and' to separate them. Put Diana as first reviewer!
\professor{Prof. Dr.-Ing. Diana Göhringer}
\chairman{Prof. Dr.-Ing. T. Mikolajick}							% Uncomment if your student is a NES student.	
%\chairman{Prof. Dr.-Ing. habil. Martin Wollschlaeger}				% Uncomment if your student is an IST student.
\taskform[pagestyle=empty]{%
Spiking Neural Networks (SNNs) are brain-inspired models that process information using discrete spike events instead of continuous values.
This makes them suitable for low-power, real-time implementations on hardware such as FPGAs.
However, efficient online learning in SNNs is still difficult to achieve in hardware.

The main goal of this thesis is to investigate efficient and flexible real-time learning in SNNs on FPGA platforms.
It focuses on existing digital SNN designs with online learning, with emphasis on the SDSP (Spike-Driven Synaptic Plasticity) algorithm.
The SDSP rule offers a hardware-friendly approach to learning through local update mechanisms.
Additionally, the thesis will explore Partial Reconfiguration (PR) techniques to enable dynamic changes of the training algorithm, while the rest of the system continues running on the FPGA.
}{% the main points of the thesis
	\item Conduct a literature review on:
    \begin{itemize}
        \item Digital implementations (FPGA or ASIC) of SDSP training algorithm.
        \item FPGA-based SNN implementations with PR-based real-time training algorithm.
    \end{itemize}
	\item Design and optimize an architectural model for the SDSP algorithm.
	\item Implement the model into RTL Code, simulate and validate it.
	\item Integrate the training module into a simplified SNN design (taking advantage of an open-source SNN framework).
	\item Make the training dynamically reconfigurable (Apply PR to the complete training block).
	\item Test it on FPGA and get (both with and without PR) the following metrics:
    \begin{itemize}
        \item Resource utilization, power consumption, latency, maximum frequency, accuracy, and for PR, reconfiguration time.
		\item (Optional): Convergence speed for learning. 
    \end{itemize}
	\item Comparison with other digital SNN architectures:
    \begin{itemize}
        \item With real-time SDSP training possibility.
        \item with PR-based real-time training possibility.
    \end{itemize}
	\item (Optional) Extend the simplified SNN to support more neuron models and/or architecture settings.
}
\end{document}
